\section{Modeling with types (the type-driven lens)}

Here is the core idea in one sentence: \textbf{a type is the set of all values
something can possibly be}, and if we describe each piece of the problem by its
type, the correct data structures and algorithms almost fall out for free.

\begin{keybox}[A note on honesty --- ``type-driven'', not ``type theory'']
What follows is \emph{type-driven design} (the practical school: strong types,
enums, \code{std::optional}, ``make illegal states unrepresentable''). It is
\emph{not} formal type theory (Curry--Howard, dependent/refinement types, proofs).
That distinction matters: in a systems take-home the pragmatic version is exactly
right, and over-claiming ``type-theoretic'' to a reviewer who knows the difference
is a needless risk. Use the vocabulary below to \emph{reason}, then implement in
plain C++17. The three ways to build types --- and this problem uses all three:
\end{keybox}

\subsection{The three atoms}

\begin{center}
\begin{tikzpicture}[font=\small,every node/.style={align=center}]
  % product
  \node[draw,rounded corners,fill=Soft,minimum width=4.6cm,minimum height=2.1cm] (p) {};
  \node[above=1pt of p.north] {\bfseries\color{NavyBlue}Product type ($A\times B$)};
  \node at (p) {``\textbf{and}''\\a value has an $A$ \emph{and} a $B$\\\footnotesize e.g.\ a \ty{Pair}, a \ty{struct}, a record};
  % sum
  \node[draw,rounded corners,fill=Soft,minimum width=4.6cm,minimum height=2.1cm,right=1.6cm of p] (s) {};
  \node[above=1pt of s.north] {\bfseries\color{Sell}Sum type ($A+B$)};
  \node at (s) {``\textbf{or}''\\a value is \emph{either} an $A$ \emph{or} a $B$\\\footnotesize e.g.\ \Buy{} \emph{or} \Sell};
  % function
  \node[draw,rounded corners,fill=Soft,minimum width=4.6cm,minimum height=2.1cm,below=1.1cm of $(p)!0.5!(s)$] (f) {};
  \node[above=1pt of f.north] {\bfseries\color{Buy}Function type ($A\to B$)};
  \node at (f) {``\textbf{turns into}''\\give it an $A$, get back a $B$\\\footnotesize e.g.\ a query, a method};
\end{tikzpicture}
\end{center}

\begin{intu}
Why the names \emph{product} and \emph{sum}? Count the values! If $A$ has 3
possible values and $B$ has 4, then a pair $A\times B$ has $3\times4=12$
possibilities (a product), while a choice $A+B$ has $3+4=7$ (a sum). The
arithmetic is literally why they're named that. This ``counting values'' habit
is the single most useful trick type theory gives an interviewer.
\end{intu}

\subsection{Typing the \texttt{Order}}

An \code{Order} is a \textbf{product type} --- it bundles several fields with
``and''. In type-theory notation:
\[
\ty{Order} \;=\; \ty{OrderId}\;\times\;\ty{SecId}\;\times\;\ty{Side}\;\times\;\ty{Qty}\;\times\;\ty{User}\;\times\;\ty{Company}
\]
and \ty{Side} is a \textbf{sum type} with exactly two inhabitants:
\[
\ty{Side} \;=\; \Buy \;+\; \Sell.
\]

\begin{pitfall}
The starter code stores \code{side} as a \code{std::string}. Type-theoretically
that is \emph{wrong}: a string has billions of possible values (\code{"Buy"},
\code{"buy"}, \code{"BUY"}, \code{"xyz"}\dots), but \ty{Side} has only
\textbf{two}. The gap between ``what the type allows'' and ``what is actually
valid'' is exactly where bugs live. In the interview, \emph{say} you would
normalize it to a 2-value \code{enum} internally (\code{Buy}/\code{Sell}) --- that
turns a whole class of invalid states into \emph{unrepresentable} states. This is
the famous principle: \emph{make illegal states unrepresentable}.
\end{pitfall}

\subsection{Typing the operations --- their signatures tell you the plan}

Read each method as a function type. The \emph{shape} of the arrow tells you what
data structure you need to make it fast.

\begin{center}
\renewcommand{\arraystretch}{1.5}
\begin{tabular}{@{}p{5.9cm}p{8.3cm}@{}}
\toprule
\textbf{Type signature} & \textbf{What the shape demands} \\
\midrule
$\code{add}:\ \ty{Order}\to\ty{Unit}$ &
returns nothing (\ty{Unit} = the 1-value type, ``void''); it's a \emph{mutation}. Must update every index. \\
$\code{cancel}:\ \ty{OrderId}\to\ty{Unit}$ &
we're handed only an \emph{id}. So we need a map \ty{OrderId}$\to$\ty{Order} to find the rest. \\
$\code{cancelForUser}:\ \ty{User}\to\ty{Unit}$ &
handed a \emph{user}. So we need \ty{User}$\to$\ty{Set of OrderId} to avoid scanning everything. \\
$\code{cancelForSecMinQty}:\ \ty{SecId}\times\ty{Qty}\to\ty{Unit}$ &
handed a \emph{security}. So we need \ty{SecId}$\to$\ty{Set of OrderId} too. \\
$\code{matchSize}:\ \ty{SecId}\to\ty{Qty}$ &
the money method: a \ty{SecId} in, a single number out. This is a \emph{fold} (aggregate) --- see \S3. \\
$\code{getAll}:\ \ty{Unit}\to\ty{List}\,\ty{Order}$ &
dump: needs a way to enumerate every live order. \\
\bottomrule
\end{tabular}
\end{center}

\begin{keybox}[The punchline of this section]
Every method that takes a key of type \ty{K} and must act on all orders with that
key is screaming for a map \ty{K}$\to$\ty{Set of OrderId}. Three different keys
(\ty{OrderId}, \ty{User}, \ty{SecId}) appear as inputs, so we will keep
\textbf{three indexes}. \emph{The type signatures alone told us how many indexes
to build.}
\end{keybox}

\subsection{The cache itself as a type}

The whole cache is a product of one \emph{source of truth} plus several
\emph{derived indexes} plus \emph{pre-computed aggregates}:
\[
\ty{OrderCache} \;=\;
\underbrace{(\ty{OrderId}\!\to\!\ty{Order})}_{\text{truth}}\;\times\;
\underbrace{(\ty{User}\!\to\!\mathcal{P}\ty{OrderId})}_{\text{index}}\;\times\;
\underbrace{(\ty{SecId}\!\to\!\mathcal{P}\ty{OrderId})}_{\text{index}}\;\times\;
\underbrace{(\ty{SecId}\!\to\!\ty{Aggregate})}_{\text{pre-sum}}
\]
where $\mathcal{P}$ means ``powerset'' = ``a set of''. The last factor is the
clever one that makes \code{matchSize} fast; we derive its exact contents next.

\begin{say}
``I'd model \code{Order} as a product type and \code{Side} as a two-value sum,
normalizing the string to an enum so illegal sides are unrepresentable. Each
operation's signature is a function type, and the key it takes tells me which
index it needs --- that's why I keep one map from each of orderId, user, and
securityId. The matching query is a fold over a per-security aggregate I maintain
incrementally.''
\end{say}
